\chapter{}

Domenico non riprese mai moglie, quantunque vi riflettesse sovente, grattandosi forte con le unghie il mento poco rasato, stringendo la pelle della gola e poi battendo le nocche su qualche cosa, ma senza farsi male. Lo annunciava con veemenza, di proposito, dopo ogni sua arrabbiatura. E credendo che Pietro si sarebbe dato agli interessi, per non trovarsi in casa una matrigna, gli diceva:

— Ora toccherebbe a te! Ma tu, imbecille, fai il socialista! Non ti vergogni?

Comprava un cappello all’anno, portandolo tutti i giorni; finchè la tesa, che si adagiava su gli orecchi, rovesciandoli più giù, non fosse untuosa. Gli piaceva di tenere la camicia almeno per due settimane; e bestemmiava quando doveva decidersi a rifarsele nuove. L’istinto di conservarsi nella condizione guadagnata lo costringeva anche ad inutili economie; che, del resto, faceva notare agli altri; anzi, volendo che fossero apprezzate, diceva, ed era vero:

— Io sono un galantuomo: ho fatto i denari con il mio sudore; e me li voglio mantenere.

In una ciotola di legno, teneva, insieme con le monete di rame, per superstizione, una medaglietta trovata mentre gli assalariati vangavano. Per guardarla meglio, il che gli succedeva tutte le volte che gli veniva in mano, mettevasi gli occhiali.

La medaglietta gli piaceva, perchè con le unghie riusciva a grattare il metallo; che, allora, pareva nuovo. Quando gli avevano portato gli occhiali, dopo averglieli cercati da per tutto, sedeva, li puliva con il suo fazzoletto rosso, puzzolente di lezzo:

— Non la vedo bene!

E usciva fuori, per farla esaminare prima al droghiere, poi al mercante e al barbiere; che erano i suoi amici più vicini.

Ma nè meno loro, naturalmente, sapevano che medaglietta fosse.

Talvolta si appoggiava, senza cappello, all’uscio della bottega; salutando anche chi conosceva a pena.

D’estate, vi si faceva portare una sedia; sonnecchiando, finchè qualcuno, che passava, non lo destasse con un colpo sopra la coscia. Allora si risentiva, dicendo:

— Mi ero addormentato un poco.

E, per levarsi il sonno, andava a dare qualche ordine.

Durante la giornata, inghiottiva tutte le frutta trapassate; e diceva al cuoco, i cui capelli neri toccavano quasi le ciglia:

— Portami un tegame!

Assaggiava e rimandava via il cuoco, spingendolo sul braccio:

— Ci hai messo poco pepe. Quando imparerai a fare da te?

Il rimproverato restava male ed alzava a poco a poco una spalla.

— Portami quell’altro tegame, ora.

Quegli obbediva, restando poi dritto a guardarlo; con una mano sopra la tavola.

Domenico non aspettava di aver ingoiato il boccone, per gridargli:

— Hai fatto bruciare l’aglio.

Si puliva i baffi, sdruciandoseli con il tovagliolo; e concludeva:

— Bisognerà che in cucina non ti lasci più solo o ti mandi via. Degli uomini non ne nascono più.

Ogni mattina mangiava di quel che c’era rimasto il giorno innanzi in fondo ai recipienti della dispensa.

Ma del vino ne beveva quasi un fiasco; e ruttava sopra il fazzoletto, volgendosi verso il muro. I sapori lo esaltavano, lo facevano loquace; e fuori della cucina gli pareva di perder tempo, a meno che non fosse a Poggio a’ Meli.
